% !TEX root = ../test_multifile.tex
\section{Chapter Two: Methods and Analysis}

This chapter presents the methods used in our analysis. We draw on the
foundational material from Chapter One and extend it with computational
techniques suitable for practical applications.

The primary tool is the discrete Fourier transform, defined for a sequence
$\{x_k\}$ of length $N$ as:
$$X_j = \sum_{k=0}^{N-1} x_k \cdot e^{-2\pi i jk / N}$$

Computing this directly requires $O(N^2)$ operations, but the Fast Fourier
Transform algorithm reduces this to $O(N \log N)$, making spectral analysis
feasible for large datasets. We implemented the Cooley-Tukey variant, which
recursively splits the problem into smaller subproblems.

Our experimental setup consisted of three benchmark datasets. The first
dataset contains synthetic signals with known frequency components, allowing
us to validate correctness. The second dataset comprises real-world audio
samples from field recordings. The third dataset is a collection of
financial time series spanning twenty years of daily observations.

For each dataset, we measured both accuracy and runtime performance. The
results confirm that our implementation matches the theoretical complexity
bounds while maintaining numerical stability across all input sizes tested.
A formula in the sub-file: $a^2 + b^2 = c^2 + zzz^3$
